\section{Cohort, instrument and denominators}
The quantitative reference cohort is the historical 423 records. The preserved accounting from 490 to 423 includes author-directed exclusions for test or non-official records and coded restrictions. It is not a response rate, a new eligibility check, or a verified count of distinct households. Do not reinstate manually excluded records or classify all 67 exclusions as tests. Referral and household links were not systematically recorded.

The five submitted form versions are 2512090157, 2512091719, 2512111044, 2512121811 and 2602120658. Unresolved versions stay unresolved. Early residence code 4 means Other; later code 4 means Soacha. Meta is an additional later option and does not identify a municipality. Use recorded dates and distinct recall windows, not inferred interview sites, exposure or migration tenure.

Keep route (open, closed or unresolved) distinct from response state: ordinary valid answer, structural skip, unavailable item or option, applicable nonresponse, item-specific sentinel, or unresolved state. A recorded-answer base may include an offered sentinel; numeric summaries use scorable answers. Multiple-response shares use the relevant valid-parent base and versions offering the option. An unoffered option, an undefined share, a withheld cell and an observed zero are different. A diagnostic flag does not alter the cohort.

\section{Outcomes and focal fields}
\begin{description}
\item[Remittance disposition (q6\_13).] Five original choices and a binary collapse: all cash (1) versus any reported digital retention/use (2--5). Codes 2 and 4 overlap; neither the category nor the collapse measures pesos, digital spending alone or the latest receipt channel.
\item[Domestic digital participation (q6\_14).] Valid No/Yes for sending or receiving within 60 days; refusal and nonresponse unscored. q6\_15 is conditional on Yes. Use the five original bins on the valid-bin base and show six response states on the affirmative-user base. Keep uncertainty (98) separate and do not reconstruct counts.
\item[Remittance context.] q6\_2 is 12-month receipt frequency; combined code 99 is not simple nonreceipt. q6\_4 records the latest receipt channel, with Other unspecified. Each has its own recorded-answer base. q6\_18 fee visibility concerns the last operation in a separate 90-day window.
\item[IAT3.] Fixed mean of the phone access/use, internet frequency and data or balance runout components. The latter can share content with the outcome. A skipped component does not become zero.
\item[IADT2.] Equal-weight average of reported phone-transfer ability (q3\_16) and fraud/message avoidance confidence (q3\_17). IADT3 adds smartphone-gated app installation (q3\_15). Neither is an observed ability test.
\item[IAFF2.] Mean of reported own-name proof of address (q4\_5) and own-name phone (q4\_6); possible values 0, 0.5 and 1. Neither legal eligibility nor account ownership is verified.
\end{description}
Retain separate item families for experiences, safety, conditional recourse, household decisions, training, support and barriers. Do not pool these into new scores or treat shared items as independent corroboration.

\section{Conditional models and uncertainty}
The E.4 working logit models share IAT3, IADT2, IAFF2, age, grouped education (completed secondary as reference) and grouped recorded residence (Bogotá as reference). R0 and D0 use remittance and domestic outcomes respectively. R1 and D1 add occupation (employee/worker as reference) and original first-arrival calendar year, not tenure. R0C refits the remittance core on the shared outcome-observed base; equality to the domestic core base is checked before omitting a duplicate domestic fit. The residual residence group is not a city or cluster.

Use each model's intended mask and estimation sample separately. Report supports, omissions, convergence and availability without inventing a refit. Changes across models can reflect covariates and composition. The fixed contrasts average predictions at the lower and higher focal values, keeping other observed covariates. IAT3 and IADT2 compare 0.50 with 0.75; IAFF2 compares 0.50 with 1.00. Contrasts are in percentage points; endpoint predictions are in probability units.

Pointwise 95\% delta-method intervals use the working robust covariance assuming independent records. The six separate coefficient-null Wald tests form one fixed Holm family. No causal, design-weighted, referral-clustered or simultaneous inference is implied.
